\pdfbookmark[1]{Abstract}{Abstract}
\chapter*{Abstract}

A battery of null results is worth only the sensitivity of the battery. The digits of pi pass every
test assembled here against SHA-256, uniform bit shares, flat autocorrelation, full algebraic degree,
no co-variation and incompressibility, and the digits of pi have a generating program of about a
kilobyte. What the battery establishes about SHA-256 is therefore a set of bounds with domains: bits
unbiased to about four parts in a million over \(550\) billion evaluations, no linear approximation
surviving out of sample, and full degree by round eight. Beside those bounds sit a real signal in the
block header, where the version-rolling bits give away \(1.82\) bits of collision entropy, and a
support measurement that needs no statistics: at round seven, \(214\) of \(256\) output positions
are forbidden to a single input bit.

Fifteen failure modes of the apparatus are listed with the subtraction that removes each. Every
instance of every mode makes a result look better than the truth. Five have a mechanical detector,
five a partial one and five none. The rules that follow are to read the interior where it is
readable, run a positive control, draw a bar by resampling and never derive it, match the null to the
nuisance parameters, normalize until the baseline is flat, measure at two widths, and suspect the
instrument before the object.

Seven transforms carry the work, and each pays only where the object has fewer coefficients than
points. A hash is built to have as many, and no transform here shortens one.
\cleardoublepage
